\documentclass{article}
\usepackage[en]{homework}

\config{Algebra\,-\,0}{2026 Spring}{4}

\begin{document}

\begin{problem}[3E-1]
	Suppose $T$ is a function from $V$ to $W$. The graph of $T$ is the subset of
	$V \times W$ defined by
	\[
		\text{graph of }T
		= \{(v, Tv) \in V \times W : v \in V\}.
	\]
	Prove that $T$ is a linear map if and only if the graph of $T$ is a subspace
	of $V \times W$.
\end{problem}

\begin{proof}
	Denote the graph of \( T \) by \( G\subset V\times W \) . On the one hand, if \( T \) is a linear map, then for and \( (v_1,Tv_1)\in G \) and  \( (v_2,Tv_2)\in G \) , and  \( \lambda\in\F \) , we have
	\[ (v_1,Tv_1)+(v_2,Tv_2)=(v_1+v_2,T(v_1+v_2))\in G,\ \lambda(v_1,Tv_1)=(\lambda v_1,T(\lambda v_1))\in G. \] 
	Hence, \( G \) is a subspace of \( V\times W \) .\par
	On the other hand, if \( G \) is a subspace of \( V\times W \) , then for any \( v_1,v_2\in V \) and \( \lambda\in\F \) , we have
	\[ (v_1,Tv_1)+(v_2,Tv_2)=(v_1+v_2,Tv_1+Tv_2)\in G,\ \lambda(v_1,Tv_1)=(\lambda v_1,\lambda Tv_1)\in G, \]
	which implies that \( T(v_1+v_2)=Tv_1+Tv_2 \) and \( T(\lambda v_1)=\lambda Tv_1 \) . Hence, \( T \) is a linear map.
\end{proof}

\begin{problem}[3E-9]
	Prove that a nonempty subset $A$ of $V$ is a translate of some subspace of
	$V$ if and only if
	\[
		\lambda v + (1 - \lambda)w \in A
	\]
	for all $v, w \in A$ and all $\lambda \in \mathbf{F}$.
\end{problem}

\begin{proof}
	On the one hand, if \( A \) is a translate of \( U \) , then there exists \( t \in V \) such that \( A = t + U \) . For any \( v, w \in A \) and \( \lambda \in \mathbf{F} \) , we have \( v = t + u_1 \) and \( w = t + u_2 \) for some \( u_1, u_2 \in U \) . Then
	\[ \lambda v+ (1-\lambda)w= \lambda(t+u_1)+(1-\lambda)(t+u_2)=t+(\lambda u_1+ (1-\lambda) u_2)\in t+U. \] \par
	On the other hand, if \( \lambda v + (1 - \lambda)w \in A \) for all \( v, w \in A \) and all \( \lambda \in \F \) , We will show that there exists  \( t\in V \) and \( U\subset V \) such that \( A=t+U \) . \par
	Pick \( a_0\in A \) and  \( a_0=t+u_0 \) . Denote
	\[ U=\left\{a-a_0:a\in A\right\}. \] \par
	It suffices to prove \( U \) is a subspace of \( V \) . 
	\begin{enumerate}[label=(\alph*)]
		\item U is closed under addition. That's because for \( u_1,u_2 \) in \( U \) , set \( a_1=a_0+u_1, a_2=a_0+u_2 \) . Then \( a_1/2+a_2/2=a_0+(u_1+u_2)/2\in A \) , hence \( 2(a_1/2+a_2/2)+(1-2)a_0=a_0+(u_1+u_2)\in A \) , which implies \( u_1+u_2\in U \) . 
		\item U is closed under scalar multiplication. That's because for \( u\in U \) and \( \lambda\in\F \) , set \( a=a_0+u \) . Then \( \lambda a+(1-\lambda)a_0=a_0+\lambda u\in A \) , which implies \( \lambda u\in U \) .
	\end{enumerate}
\end{proof}

\begin{problem}[3E-14]
	Suppose $U$ and $W$ are subspaces of $V$ and $V = U \oplus W$. Suppose
	$w_1, \cdots, w_m$ is a basis of $W$. Prove that
	$w_1 + U, \cdots, w_m + U$ is a basis of $V/U$.
\end{problem}

\begin{proof}
	First we will show that \( w_1+U, \cdots, w_m+U \) is linearly independent. Suppose \( \lambda_1(w_1+U)+\cdots+\lambda_m(w_m+U)=0+U \) for some \( \lambda_1,\cdots,\lambda_m\in\F \) . Then \( \lambda_1w_1+\cdots+\lambda_mw_m\in U \) . Since \( U\cap W=\{0\} \) , we have \( \lambda_1w_1+\cdots+\lambda_mw_m=0 \) , which implies \( \lambda_1=\cdots=\lambda_m=0 \) . Hence, \( w_1+U, \cdots, w_m+U \) is linearly independent.\par
	Next we will show that \( w_1+U, \cdots, w_m+U \) spans \( V/U \) . For any \( v+U\in V/U \) , we can write \( v=u+w \) for some \( u\in U \) and \( w\in W \) . Since \( w_1, \cdots, w_m \) is a basis of \( W \) , we can write \( w=\lambda_1w_1+\cdots+\lambda_mw_m \) for some \( \lambda_1,\cdots,\lambda_m\in\F \) . Then
	\[ v+U=u+w+U=\lambda_1(w_1+U)+\cdots+\lambda_m(w_m+U). \]
	Hence, \( w_1+U, \cdots, w_m+U \) spans \( V/U \) .\par
	Combining the above two paragraphs, we conclude that \( w_1+U, \cdots, w_m+U \) is a basis of \( V/U \) .
\end{proof}

\begin{problem}[3E-15]
	Suppose $U$ is a subspace of $V$ and $v_1 + U, \cdots, v_m + U$ is a basis
	of $V/U$ and $u_1, \cdots, u_n$ is a basis of $U$. Prove that
	$v_1, \cdots, v_m, u_1, \cdots, u_n$ is a basis of $V$.
\end{problem}

\begin{proof}
	First we will show that \( v_1, \cdots, v_m, u_1, \cdots, u_n \) is linearly independent. Suppose \( \lambda_1v_1+\cdots+\lambda_mv_m+\mu_1u_1+\cdots+\mu_nu_n=0 \) for some \( \lambda_1,\cdots,\lambda_m,\mu_1,\cdots,\mu_n\in\F \) . Then \( \lambda_1(v_1+U)+\cdots+\lambda_m(v_m+U)=0+U \) , which implies \( \lambda_1=\cdots=\lambda_m=0 \) . Hence, \( \mu_1u_1+\cdots+\mu_nu_n=0 \) , which implies \( \mu_1=\cdots=\mu_n=0 \) . Therefore, \( v_1, \cdots, v_m, u_1, \cdots, u_n \) is linearly independent.\par
	Next we will show that \( v_1, \cdots, v_m, u_1, \cdots, u_n \) spans \( V \) . Since \( v_1+U, \cdots, v_m+U \) is a basis of \( V/U \) , we can write \( v+U=\lambda_1(v_1+U)+\cdots+\lambda_m(v_m+U) \) for some \( \lambda_1,\cdots,\lambda_m\in\F \) . Then
	\[ v=\lambda_1v_1+\cdots+\lambda_mv_m+u \] for some \( u\in U \) . Since \( u_1, \cdots, u_n \) is a basis of \( U \) , we can write \( u=\mu_1u_1+\cdots+\mu_nu_n \) for some \( \mu_1,\cdots,\mu_n\in\F \) . Then
	\[ v=\lambda_1v_1+\cdots+\lambda_mv_m+\mu_1u_1+\cdots+\mu_nu_n. \]
	Hence, \( v_1, \cdots, v_m, u_1, \cdots, u_n \) spans \( V \) .\par
	Combining the above two paragraphs, we conclude that \( v_1, \cdots, v_m, u_1, \cdots, u_n \) is a basis of \( V \) .
\end{proof}

\begin{problem}[3E-18]
	Suppose that $U$ is a subspace of $V$ such that $V/U$ is finite-dimensional.
	\begin{enumerate}[label=(\alph*)]
		\item Show that if $W$ is a finite-dimensional subspace of $V$ and $V = U + W$, then $\dim W \geq \dim V/U$.
		\item Prove that there exists a finite-dimensional subspace $W$ of $V$ such that $\dim W = \dim V/U$ and $V = U \oplus W$.
	\end{enumerate}
\end{problem}

\begin{proof}
	(a)  Suppose \( W \) is a finite-dimensional subspace of \( V \) and \( V = U + W \) . Then the map \( T:W\to V/U \) defined by \( Tw=w+U \) is a surjective linear map. Hence, we have
	\[ \dim W=\dim\ker T+\dim V/U\geq \dim V/U. \] \par
	(b)  Since \( V/U \) is finite-dimensional, we can pick a basis \( v_1+U, \cdots, v_m+U \) of \( V/U \) . Denote \( W=\sspan\{v_1,\cdots,v_m\} \) . We will show that \( V=U\oplus W \) and \( \dim W=\dim V/U \) .\par
	First we will show that \( V=U+W \) . For any \( v\in V \) , we can write \( v+U=\lambda_1(v_1+U)+\cdots+\lambda_m(v_m+U) \) for some \( \lambda_1,\cdots,\lambda_m\in\F \) . Then
	\[ v=\lambda_1v_1+\cdots+\lambda_mv_m+u \] for some \( u\in U \) . Hence, \( V=U+W \) .\par
	Next we will show that \( U\cap W=\{0\} \) . Suppose \( u\in U\cap W \) . Then \( u=\lambda_1v_1+\cdots+\lambda_mv_m \) for some \( \lambda_1,\cdots,\lambda_m\in\F \) . Hence, we have
	\[ \lambda_1(v_1+U)+\cdots+\lambda_m(v_m+U)=0+U, \] which implies \( \lambda_1=\cdots=\lambda_m=0 \) . Therefore, \( u=0 \) and \( U\cap W=\{0\} \) .\par
	Combining the above two paragraphs, we conclude that \( V=U\oplus W \) .\par
	Finally, since \( v_1+U, \cdots, v_m+U \) is a basis of \( V/U \) , we have \( \dim W=\dim\sspan\{v_1,\cdots,v_m\}=m=\dim V/U \) .
\end{proof}

\begin{problem}[3F-5]
	Suppose $T \in \mathcal{L}(V, W)$ and $w_1, \cdots, w_m$ is a basis of $\im T$. Hence for each $v \in V$, there exist unique numbers $\varphi_1(v), \cdots, \varphi_m(v)$ such that
	\[
		Tv = \varphi_1(v)w_1 + \cdots + \varphi_m(v)w_m,
	\]
	thus defining functions $\varphi_1, \cdots, \varphi_m$ from $V$ to $\F$. Show that each of the functions $\varphi_1, \cdots, \varphi_m$ is a linear functional on $V$.
\end{problem}

\begin{proof}
	Without loss of generality, we only need to show that \( \varphi_1 \) is a linear functional on \( V \) . For any \( v_1,v_2\in V \) and \( \lambda\in\F \) , we have
	\[ \varphi_1(v_1+v_2)w_1+\cdots=v_1+v_2=(\varphi_1(v_1)+\varphi_1(v_2))w_1+\cdots\implies \varphi_1(v_1+v_2)=\varphi_1(v_1)+\varphi_1(v_2), \] 
	\[ \lambda\varphi_1(v_1)w_1+\cdots=\lambda v_1=\varphi_1(\lambda v_1)w_1+\cdots\implies \varphi_1(\lambda v_1)=\lambda\varphi_1(v_1). \] 
\end{proof}

\begin{problem}[3F-8]
	Suppose $v_1, \cdots, v_n$ is a basis of $V$ and $\varphi_1, \cdots, \varphi_n$
	is the dual basis of $V'$. Define $\Gamma : V \to \F^n$ and
	$\Lambda : \F^n \to V$ by
	\[
		\Gamma(v) = (\varphi_1(v), \cdots, \varphi_n(v))
	\]
	and
	\[
		\Lambda(a_1, \cdots, a_n) = a_1v_1 + \cdots + a_nv_n.
	\]
	Explain why $\Gamma$ and $\Lambda$ are inverses of each other.
\end{problem}

\begin{proof}
	\[ \Lambda\Gamma(v)=\Lambda(\varphi_1(v),\cdots\varphi_n(v))=\varphi_1(v)v_1+\cdots+\varphi_n(v)v_n=v, \]
	\[ \Gamma\Lambda(a_1,\cdots,a_n)=\Gamma(a_1v_1+\cdots+a_nv_n)=(\varphi_1(a_1v_1+\cdots+a_nv_n),\cdots,\varphi_n(a_1v_1+\cdots+a_nv_n))=(a_1,\cdots,a_n). \]
\end{proof}

\begin{problem}[3F-10]
	Suppose $m$ is a positive integer.
	\begin{enumerate}[label=(\alph*)]
		\item Show that $1, x - 5, \cdots, (x - 5)^m$ is a basis of
		      $\mathcal{P}_m(\mathbb{R})$.
		\item What is the dual basis of the basis in (a)?
	\end{enumerate}
\end{problem}

\begin{solution}
	(a) For any \( p\in \mathcal{P}_m(\R) \) , we can write uniquely
	\[ p(x)=\sum_{k=0}^{m} \frac{p^{(k)}(5)}{k!} (x-5)^k, \] 
	which is not hard to verify.\par
	(b) From the formula in (a), we can see that the dual basis of \( 1, x - 5, \cdots, (x - 5)^m \) is
	\[ \varphi_k(p)=\frac{p^{(k)}(5)}{k!},\ 0\le k\le m. \] 
\end{solution}

\begin{problem}[3F-11]
	Suppose $v_1, \cdots, v_n$ is a basis of $V$ and $\varphi_1, \cdots, \varphi_n$
	is the corresponding dual basis of $V'$. Suppose $\psi \in V'$. Prove that
	\[
		\psi = \psi(v_1)\varphi_1 + \cdots + \psi(v_n)\varphi_n.
	\]
\end{problem}

\begin{proof}
	For any \( v\in V \) , we can write \( v=\varphi_1(v)v_1+\cdots+\varphi_n(v)v_n \) . Applying \( \psi \) to both sides, we get
	\[ \psi(v) = \psi(\varphi_1(v)v_1+\cdots+\varphi_n(v)v_n) = \varphi_1(v)\psi(v_1)+\cdots+\varphi_n(v)\psi(v_n), \] 
	which is exactly we want to prove.
\end{proof}



\end{document}